
\makeatletter
\let\@oldaddcontentsline\addcontentsline
\renewcommand{\addcontentsline}[3]{}
\section*{附录}
\let\addcontentsline\@oldaddcontentsline
\makeatother

\begin{table}[H]
  \centering
  \caption{附件说明表}
  \label{tab:appendix}
  \renewcommand{\arraystretch}{1.2}
  \begin{tabularx}{\textwidth}{LX}
    \toprule
    文件 & 说明 \\
    \midrule
    问题一\_单颗粒导通判断.py & 问题一求解代码（单颗粒导通判定） \\
    问题二\_多颗粒蒙特卡洛.py & 问题二求解代码（多颗粒蒙特卡洛模拟） \\
    问题三\_最小填充率搜索.py & 问题三求解代码（二分搜索最小填充率） \\
    问题四\_双金属成本优化.py & 问题四求解代码（双金属成本优化） \\
    core\_geometry.py & 核心几何算法模块（距离计算、连通判定） \\
    附件.xlsx（附件1） & 12个金属A颗粒的端点坐标 \\
    附件.xlsx（附件2） & 49个金属A颗粒的端点坐标 \\
    附件.xlsx（附件3） & 535个金属A颗粒的端点坐标 \\
    求解计划.md & 完整求解计划文档 \\
    问题一判定结果.csv & 12个颗粒的导通判定详细结果 \\
    问题二结果汇总.csv & 各填充率蒙特卡洛模拟结果 \\
    问题三搜索结果.csv & 二分搜索过程与最小填充率结果 \\
    问题四最优配置.csv & 双金属成本优化最优方案 \\
    \bottomrule
  \end{tabularx}
\end{table}
